\section{Complexity, effective algebra, and finite geometric tests}
\label{sec:complexity}

Here $A_D$ denotes the rational Dirichlet-GRH input designated by the task as
Alp\"oge's theorem. The letter $M$ denotes Masullo's all-number-field Hecke-GRH
claim, and every use of $M$ remains conditional. The input size of an arithmetic
problem includes the circuits, explicitly listed variables, rational constants
in binary, and any promised field representation. A geometric object without
an encoding has no associated complexity bound.

\subsection{Classical consequences and inherited results}

The following catalogue entries preserve the original finding identifiers of
the earlier runs. They are not advertised as new findings of this session.
An exact locator to an earlier proof is distinguished from a fresh reading of
that proof. In particular, the earlier complete file \texttt{CX\_GRH.tex} is
not in tonight's attachment: its \texttt{IMPLICATIONS}, \texttt{LEDGER},
\texttt{QUESTIONS}, and targeted intake are available. Statements that depend
on its elaborate new finite-field constructions are carried from those records;
this session does not manufacture an independent verification of the absent
proof text.

\paragraph{CX-ALG-001 and CX-ALG-002: primality.}
Bach, \emph{Explicit bounds for primality testing and related problems}
(1990), Theorem 2, p.~372, and its application on p.~373 give the cutoff
$2(\log N)^2$ for a compositeness witness. Schoof, \emph{Four primality testing
algorithms} (2008), Theorem 2.2, pp.~105--106, states and proves the Miller test
with this cutoff using only quadratic Dirichlet characters. Therefore $A_D$
discharges its analytic input. Writing $N-1=2^s m$ with $m$ odd, test the usual
conditions $a^m=1$ or $a^{2^i m}=-1$ modulo $N$ for every base through the
cutoff, also checking nontrivial gcds. Even inputs are handled separately.
There are $O((\log N)^2)$ bases and each modular exponentiation takes a
polynomial number of bit operations. With $b=\lceil\log_2N\rceil$ and
multiplication cost $\mathcal M(b)$, the elementary implementation uses
$O(b^3\mathcal M(b))$ bit operations. This is the transfer already recorded
as CX\_GRH CX-N01 and ALPOGE\_GCT C-GRH-02.

The statement $\mathrm{PRIMES}\in\mathrm P$ was already unconditional:
Agrawal--Kayal--Saxena, \emph{PRIMES is in P}, Annals 160 (2004), Theorems 4.1
and 5.1, with their 2019 erratum. The new input validates this particular
short-witness test; it does not first place primality in $\mathrm P$.
Geometrically these are computations in $\mathbf G_m$ over $\mathbf Z/N\mathbf Z$
and, for AKS, finite cyclotomic algebras. The archival companion is cote 28,
scan 3 (group-action axioms), as mapped in the CX\_GRH return, not freshly
inspected from an archive scan in this session.

\paragraph{CX-ALG-003--CX-ALG-004: generators and nonresidues.}
Bach's Theorem 3, p.~372, gives units below $3(\log N)^2$ generating
$(\mathbf Z/N\mathbf Z)^\times$. Applied to the proper subgroup of $r$th
powers in $\mathbf F_p^\times$, it gives a polynomial search for a specified
prime $r\mid p-1$. Membership is tested by exponentiation to $(p-1)/r$;
the other factors of $p-1$ need not be known. This is CX\_GRH CX-N02.
The original Bach PDF was not freshly recovered here (the publisher returned
403); its exact locator is carried from the preceding checked run. Schoof's
primary source independently supplies the quadratic special case used above.
In the Kummer sequence this is finding a nonzero class in
$\mathbf F_p^\times/(\mathbf F_p^\times)^r$, or a nontrivial
$\boldsymbol\mu_r$-torsor. The established companion locator is SGA 4 IX 3.2.

\paragraph{CX-ALG-006--CX-ALG-008: primitive roots and extension fields.}
CX\_GRH CX-N03 constructs a certified generator of $\mathbf F_p^\times$ when
the full factorization of $p-1$ is supplied, by constructing each primary
component and multiplying them. Shoup, \emph{Searching for Primitive Roots
in Finite Fields} (1992), author version Theorem 3, p.~4, and Section 6,
p.~13, gives a least primitive root of order
$O(r^4(\log r+1)^4(\log p)^2)$, where $r=\omega(p-1)$,
and hence $O((\log p)^6)$ under the
rational Dirichlet hypothesis. This is a guaranteed short candidate list
without a supplied factorization, as CX\_GRH CX-N04 explicitly records.
A candidate list does not identify or certify a particular primitive root.

The stronger extension-field construction CX\_GRH CX-N13 supplies a specified
$r$th nonresidue in a supplied $\mathbf F_{p^m}$ in
$\operatorname{poly}(m\log p)$ time when $r\mid p-1$ is prime. Its quadratic
case with even $m$ already has an unconditional construction. The theorem is
inherited from the full targeted check in
\texttt{OURS\_INTAKE\_CX\_GRH\_20261121.md}, Section 1. A fresh source audit
found a nonvanishing gap in one cited auxiliary formula; CX-ALG-016 gives
the exact counterexample, and CX-ALG-017--CX-ALG-018 give a fully checked
repair and re-proof of the needed construction. The defining
irreducible polynomial is part of the input. These are computations with
Kummer covers inside an explicitly represented finite \`etale algebra.

\paragraph{CX-ALG-009--CX-ALG-010: polynomial factorization over finite fields.}
The inherited theorem CX\_GRH CX-N15 has bound
\[
                 (d^{\log d}\log q)^{O(1)}
\]
for a degree-$d$ polynomial over a supplied $\mathbf F_q$. Under $M$ its
analytic input is satisfied by $\zeta_{\mathbf Q}$ and the particular
$\zeta_{E_\nu}$ used in the construction, where
$E_\nu=\mathbf Q(\zeta_{2\nu},p^{1/\nu})$ and $\nu\leq d^2$.
Its exact component locators are Adleman--Lenstra (1986), Theorem 1, p.~350,
Algorithm A, pp.~350--351; Bach--von zur Gathen--Lenstra (2001), Theorem 2,
p.~7; and the 2008 Ivanyos--Karpinski--R\'onyai--Saxena version, Main Theorem,
p.~3, Lemma 4.2, p.~14, Proposition 4.6, pp.~16--17.
The fresh primary published IKRS article confirms the quasipolynomial
factorization-versus-automorphism distinction (Theorems 1.1 and 1.3), and
its unconditional noncommutative zero-divisor result (Theorem 1.7).

The displayed bound remains quasipolynomial in $d$. The residual $A_D$-only
problem in CX\_GRH CX-U02 is the high-valuation branch
$p>d$, $r\nmid p(p-1)$ and
$v_r(p^{\operatorname{ord}_r(p)}-1)>1$: construct the maximal coefficient
Teichm\"uller unit needed by that descent, or replace the descent. No new
solution of that gate is asserted. Geometrically factorization computes
primitive idempotents, hence connected components, of a finite \`etale
algebra (SGA 1 V Section 7). It is a different computational problem from
factoring an integer with unknown prime divisors.

\paragraph{CX-ALG-011--CX-ALG-013: arithmetic Picard groups.}
Bach (1990), Theorem 4, p.~375, and its corollary, p.~376, give the familiar
$12(\log|D_K|)^2$ class-group generator bound, using class-character Hecke
$L$-functions over $K$. Thus $M$ supplies the premise. This is CX\_GRH
CX-N07 and ALPOGE\_BRIDGES CC-F03/F03b. If the Hilbert class field is
abelian over $\mathbf Q$, $A_D$ suffices by Dirichlet factorization; the
weaker condition that $K/\mathbf Q$ be abelian does not establish this.
A factor base is not the full lattice of relations between its ideal
classes. Relation collection, saturation, regulator information and their
running times remain separate inputs in CX\_GRH CX-N08 and CX-U03.
In Grothendieck's language the class group is
$\operatorname{Pic}(\operatorname{Spec}\mathcal O_K)$ and principal divisors
supply the relation lattice. The established companion locators are SGA 4
IX Theorem 3.3, SGA 6 V 3.10 and VI 3.2--3.3, and cote 2 scan 5 as recorded
by CX\_GRH.

\paragraph{CX-ALG-019: a complete class-group algorithm retains randomness.}
For a supplied negative fundamental discriminant $-d$, Biasse--Erukulangara,
\emph{A proof of the conjectured run time of the Hafner--McCurley class
group algorithm}, Main Theorem, author PDF p.~4, Algorithm 4, p.~15,
and Corollary 1, p.~16, give a Las Vegas computation of
$\operatorname{Cl}(\mathcal O_{\mathbf Q(\sqrt{-d})})$ with expected bound
\[
 \exp\!\left(\left(\frac3{\sqrt8}+o(1)\right)
                         \sqrt{\log d\log\log d}\right).
\]
Their Conjecture 1, p.~4, explicitly states number-field Hecke-GRH; hence
$M$ supplies the hypothesis. The proof also uses the class-group Cayley
graph expansion of Jao--Miller--Venkatesan (2009), Theorem 1.1 and
Corollary 1.3. Rational quadratic Dirichlet-GRH alone does not supply all
these class-character estimates. This is the complete algorithmic transfer
already recorded as CX\_GRH CX-N12, not a new algorithm of this session.

The source's bounded trial returns a certified answer or FAILURE, with
success probability $1-d^{-1+o(1)}$ and the displayed time bound. Repeating
on detected failure gives the stated expected Las Vegas bound. The
certification in Algorithm 4 compares the order $h'$ of the group presented
by the collected relations with $2h^*$, where $h^*\le h<2h^*$ and $h$
is the actual class number. Since $h\mid h'$ and $h'>0$, the condition
$h'<2h^*$ forces $h'=h$. Thus the restart does not accept an incorrect
relation lattice. The full earlier proof and its detailed source check are
credited to CX-N12; the primary theorem, analytic-input definition,
Algorithm 4 and Corollary 1 were freshly inspected here.
The runtime is subexponential in $\log d$, not polynomial in that input
length, and the algorithm remains randomized. Its geometric output is the
finite Picard group, with actual principal-divisor relations rather than
only a factor base. The prior companions are SGA 4 IX Theorem 3.3,
SGA 6 V 3.10 and VI 3.2--3.3, and cote 2 scan 5.

\paragraph{CX-ALG-020: small-prime walks in ray class groups.}
Jao--Miller--Venkatesan (2009), Theorem 1.1, author PDF p.~3,
Remark 1.2(a) and Corollary 1.3, p.~4, give the following exact transfer.
For degree-$n$ $K$, modulus $\mathfrak m$, and
$Q=|D_K|N\mathfrak m$, take the narrow ray class group $G$ and the
multiset of degree-one prime ideals of norm at most $x$, coprime to
$\mathfrak m$, together with their inverses. Conditional on $M$, its
Cayley adjacency eigenvalues satisfy
\[
 \lambda_1=2\operatorname{li}(x)+O(n\sqrt x\log(xQ)),\qquad
 |\lambda_\chi|=O(n\sqrt x\log(xQ))\quad(\chi\ne1).
\]
The constants here are absolute. For fixed $n$ and $B>2$,
$x\ge(\log Q)^B$, and sufficiently large $Q$, a walk of length
$t\ge C_{n,B}\log|G|/\log\log Q$ reaches any fixed subset $T$ with
probability at least $|T|/(2|G|)$. Quotients inherit the bound.
The precise hypothesis is GRH for the Hecke characters actually occurring;
for $K=\mathbf Q$, $A_D$ suffices. Random walks remain randomized.
Geometrically these are bounded-prime correspondences on arithmetic
Picard/ray-class quotients, complementing CX\_GRH CX-N12; the Picard
companions are SGA 4 IX Theorem 3.3 and cote 2 scan 5.

\paragraph{CX-ALG-022: the ordinary elliptic isogeny graph.}
The same paper's Theorem 1.5, p.~5, applies conditionally on $M$ to ordinary
elliptic curves over a fixed $\mathbf F_q$, with fixed point count and
the same geometric endomorphism order. The graph of isogenies of prime
degree less than $(\log(4q))^B$, $B>2$ fixed, has the stated expansion
bound. Theorem 1.6, pp.~6--7, transfers a discrete-logarithm solver working
on a fraction $\mu$ of that level to every curve on it with an expected
$\mu^{-1}\operatorname{poly}(\log q)$ overhead. These are conditional
classical transfers; they do not supply such a solver. The Hecke input is
for imaginary-quadratic ring-class characters, rather than an assertion
about a non-CM elliptic curve's global L-function. The geometric link is
the CM ideal-class action (Proposition 3.1, p.~11) and Deuring lifting
(Theorem 3.3, p.~12); ORDER item 20 supplies the abelian-scheme/isogeny
programme, alongside the inherited Picard comparison of cote 2 scan 5.

\paragraph{CX-HN-001--CX-HN-004: the exact Nullstellensatz input.}
Koiran's DIMACS 96-27, Section 2, Theorem 2, pp.~5--6, proves
$\mathrm{HN}\in\mathrm{AM}\subseteq\Pi_2^{\mathrm P}$ under its GRH input.
The August 8, 1997 erratum corrects the lower estimate in Theorem 8 and a
prime-product expression in Theorem 10; it explicitly preserves the main
statements. The primary author-hosted report and the entire erratum were
retrieved in this session. Section 5, Theorem 9 and Corollary 1, printed p.~12 (PDF page~13),
show that the analytic estimate is for the bounded primitive solution
field itself. A Galois closure is not required. These are the earlier
CX\_GRH CX-K01--K02; the circuit and homogeneous nonzero forms are also
stated in GCT V Theorem 2.10, p.~25.

Consequently $M$ gives the unrestricted transfer, whereas $A_D$ gives the
bounded abelian-witness promise of CX\_GRH CX-K03. All the arithmetic size
bounds in that promise matter: field degree, logarithmic discriminant and
denominators have the stated single-exponential bounds in the input length.
The universal replacement by an abelian solution field was already refuted
by $x^3-2$ in CX\_GRH CX-K04 and ALPOGE\_GCT C-GRH-03. This refutes that
specific replacement, not every possible future $A_D$-only proof.
The geometric object is a finite \`etale spread of the solution field over
an open of $\operatorname{Spec}\mathbf Z$; degree-one Frobenius orbits give
rational reductions. The companion locators are SGA 1 V Section 7 and
cote 98 scans 11--12 as carried in CX\_GRH.

\paragraph{CX-NN-001--CX-NN-004: normalization maps.}
The preceding CX\_GRH CX-K05 refines GCT V Theorem 3.4, pp.~27--28, to
$\mathrm{FP}^{\Sigma_2^{\mathrm P}}$ under $M$ for equation-given,
standard-graded algebras with explicit ambient variables and polynomial
output bits. The source proof's prefix-existence queries already have
$\Sigma_2$ form. Its unconditional polynomial-space statement remains
available. The resulting map
$\operatorname{Spec}R\to\mathbf A^{\dim R}$ is finite; the companion is
EGA II 6.1.4--6.1.5, and cote 98 scan 44 in the prior crosswalk.

The historical EXPH transfers for succinct small separating systems are
GCT V Theorems 4.10, p.~34, and 5.12, carried in CX\_GRH CX-K06. The
optimal-cardinality h.s.o.p. bounds of Theorems 4.1, p.~30, and 5.5 retain
their efficient defining-equation premise. Their output conventions cannot
be silently changed. Earlier unconditional improvements already exist:
Forbes--Shpilka, ECCC TR17-163, Theorem 7.1, p.~21, constructs the relevant
hitting set in polynomial space for unary $n,s,r$; its introductory
Theorems 1.2--1.3 record normalization consequences. Guo--Saxena--Sinhababu
(2019), Theorems 1.5--1.6, are recorded in CX\_GRH CX-K07. These are not
new consequences of tonight's analytic inputs.

\subsection{New formulations with complete proofs}

\paragraph{CX-ALG-005: deterministic prime-field square roots.}
\emph{Statement.} Under $A_D$, given an odd prime $p$ and $a\in\mathbf F_p$,
one can find all roots of $X^2-a$ or certify their absence in deterministic
time polynomial in $\log p$.

\emph{Proof.} The case $a=0$ is immediate. Euler's criterion decides whether
$a$ is a square. If it is not, return no roots. Otherwise find a quadratic
nonresidue by the deterministic search of CX-ALG-004. Apply the classical
Tonelli--Shanks algorithm with this nonresidue. More specifically, Algorithm
3.34 of Menezes--van Oorschot--Vanstone, \emph{Handbook of Applied Cryptography},
p.~100, has only one random step, its choice of the nonresidue in Step 2.
All its remaining steps are modular exponentiation, gcd inversion and a
loop of length at most $v_2(p-1)\leq\log_2p$; Fact 3.35 gives a polynomial
bit bound for those operations. Replacing Step 2 by the bounded deterministic
search preserves polynomial time and the algorithm's proved correctness.
The resulting roots are $r,-r$. Completing the square factors every quadratic
with nonzero leading coefficient; in characteristic two the separate field
case is trivial for $\mathbf F_2$. This does not provide a square-root oracle
modulo an integer whose factors are unknown.\hfill$\square$

\paragraph{CX-ALG-014: a compositeness witness need not give the usual gcd factor.}
\emph{Statement.} The rule ``take a Miller witness and compute
$\gcd(a,N)$ or $\gcd(a^m\pm1,N)$'' does not always output a nontrivial divisor.

\emph{Proof.} Take $N=35$, so $N-1=2\cdot17$, and $a=2$. Since
$2^{12}=1\pmod{35}$, one has $2^{17}=32\pmod{35}$. This is neither $1$ nor
$-1$ modulo $35$, so $2$ is a Miller witness. Yet
\[
 \gcd(2,35)=\gcd(32-1,35)=\gcd(32+1,35)=1.
\]
This is a counterexample to precisely the displayed extraction rule.
It is not a lower bound for general integer factorization and does not
rule out other algorithms.\hfill$\square$

\paragraph{CX-ALG-015: integer factorization with a supplied exponent multiple.}
\emph{Statement.} Given integers $N\ge2$ and $T\ge1$, with the promise
\[
       \lambda'(N):=\operatorname{lcm}_{p\mid N}(p-1)\mid T,
\]
the complete factorization of $N$ is computable deterministically in time
polynomial in $\log N+\log T$ under $A_D$. This is the input-length form
of Miller (1976), Lemma 5 and Theorem 3, p.~310. In particular, a supplied
value of $\varphi(N)$ or $\lambda(N)$ meets the promise. Neither value is
asserted to be available for an arbitrary $N$.

\emph{Proof of the stated input-length form.} Remove powers of two, test
primality, and detect perfect powers by integer-root computations. These
are polynomial-time operations. For a perfect power, recurse on its base
and multiply the resulting exponents. It remains to split an odd composite
with at least two distinct prime divisors. Write $T=2^s t$ with $t$ odd.
For every $a\leq\lceil3(\log N)^2\rceil$, first compute $\gcd(a,N)$;
return a proper nontrivial gcd if found, and otherwise test, for units $a$,
\[
             \gcd\bigl(a^{2^j t}-1,N\bigr),\qquad 0\leq j<s.
\]
Powers are computed modulo $N$ by repeated squaring.

To prove that some test splits $N$, put
$e_p=v_2(p-1)$ and $e=\max_{p\mid N}e_p$. If $e_p=e$ and $e_q<e$ for
two divisors, the quadratic character $a\mapsto(a/p)$ is a nontrivial
character of $(\mathbf Z/N\mathbf Z)^\times$. If all $e_p=e$, use
$a\mapsto(a/p)(a/q)$ for distinct divisors $p,q$ instead. This too is
nontrivial, by the Chinese remainder theorem. Bach's generating bound
CX-ALG-003 implies that some tested unit has character $-1$.
In the first case its reduction modulo $p$ has multiplicative order with
2-part $2^e$, whereas the 2-part modulo $q$ is smaller. In the second
case one Legendre symbol is $-1$ and the other is $1$, with the same
consequence after interchanging $p,q$. Since $t$ is a multiple of the odd
part of each $p-1$, one obtains
\[
       a^{2^{e-1}t}=-1\pmod p,\qquad
       a^{2^{e-1}t}=1\pmod q.
\]
The corresponding gcd is divisible by $q$ and not by $p$, so is proper.
The promise gives $s\ge e$, hence this test occurs in the list. This
argument uses only reductions modulo the distinct primes; repeated prime
factors in $N$ cause no problem.

Every divisor $D$ encountered in the recursion still satisfies
$\lambda'(D)\mid T$. A recursion tree has $O(\log N)$ leaves and nodes
after perfect-power compression. Each node uses polynomially many bases
and at most $\log_2T$ squarings per base, so the total bit cost is
polynomial in the displayed input length. Geometrically the known exponent
annihilates the multiplicative groups of the reduced residue fibers;
quadratic characters separate two unknown fibers and a gcd recovers their
different supports. This is a precise promise case of the integer-factoring
gate in CX-ALG-014, not an algorithm for producing the exponent promise.
\hfill$\square$

\paragraph{CX-ALG-016: the fixed-generator resolvent needs a nonzero test.}
\emph{Statement.} The literal fixed-generator formula in
Bhargava--Ivanyos--Mittal--Saxena (2017), Theorem 3.1, author PDF p.~3,
can be zero. This does not refute their nonresidue-construction conclusion;
CX-ALG-017 supplies the missing nonzero choice.

\emph{Proof.} Let $q=7$, $r=3$, $d=3$, and
$E=\mathbf F_7[X]/(X^3-2)$, with $\alpha=X\bmod(X^3-2)$.
The nonzero cubes in $\mathbf F_7$ are $1,-1$, so this cubic has no root
and is irreducible. Take the primitive cube root $\zeta=4$. Then
$\alpha^7=4\alpha$ and $\alpha^{49}=2\alpha$, whence
\[
 L=\alpha+4\alpha^7+2\alpha^{49}
   =(1+16+4)\alpha=0.
\]
Consequently $L^{(7^3-1)/3}=0$, whereas the stated right-hand side is
$\zeta^{-1}=2$. The displayed formula and its hypotheses were checked
against the rendered original PDF page, not merely its text extraction.
The invalid step is division by $L$ when inferring
$L^{q^k-1}=\zeta^{-1}$ from $L^{q^k}=\zeta^{-1}L$.
\hfill$\square$

\paragraph{CX-ALG-017: a basis-search repair for the resolvent.}
\emph{Statement.} Let $E=\mathbf F_q[X]/(f)$ be a supplied field of degree
$d=rk$, where $r$ is prime. Suppose a primitive $r$th root of unity
$\zeta\in\mathbf F_q$ is supplied, and assume $r\mid q-1$ for odd $r$,
or $4\mid q-1$ when $r=2$. Then an $r$th nonresidue in $E$, and its
nonresidue norm in $\mathbf F_q$, can be found deterministically in time
polynomial in $d\log q$. No analytic hypothesis is needed for this statement.

\emph{Proof.} Write $Q=q^k$ and define the $\mathbf F_q$-linear map
\[
       P(z)=\sum_{i=0}^{r-1}\zeta^i z^{Q^i}.
\]
As a polynomial in $z$, it is nonzero and has degree
$Q^{r-1}<Q^r=|E|$. Thus it cannot vanish at every element of $E$.
Since it is linear, it cannot vanish on every element of the power basis
$1,\alpha,\ldots,\alpha^{d-1}$. Scan this basis and choose
$\beta=P(\alpha^j)\ne0$. Cyclically shifting the sum and using
$z^{Q^r}=z$ gives $\beta^Q=\zeta^{-1}\beta$, hence
$\beta^{Q-1}=\zeta^{-1}$.

Put $S=1+Q+\cdots+Q^{r-1}$. For odd $r$, write $Q=1+rc$ and expand
modulo $r^2$ to get
\[
 S\equiv r+rc\frac{r(r-1)}2\equiv r\pmod{r^2}.
\]
For $r=2$, the hypothesis gives $S=1+Q\equiv2\pmod4$.
In both cases $S/r\equiv1\pmod r$. Therefore
\[
 \beta^{(q^d-1)/r}
   =\left(\beta^{Q-1}\right)^{S/r}=\zeta^{-1}\ne1.
\]
This proves that $\beta$ is an $r$th nonresidue. Its norm satisfies
\[
 \left(N_{E/\mathbf F_q}\beta\right)^{(q-1)/r}
      =\beta^{(q^d-1)/r}=\zeta^{-1},
\]
so is also a nonresidue. There are $d$ basis candidates and $r\le d$
terms per test; all exponent bit lengths are $O(d\log q)$.
Field arithmetic in the supplied representation and these exponentiations
therefore take polynomial time.\hfill$\square$

For the inherited CX\_GRH CX-N13, obtain $\zeta$ under $A_D$ from a
prime-field nonresidue. If odd $r\mid m$, apply the theorem with
$q=p,d=m$. If $r\nmid m$, a prime-field nonresidue $a$ remains a
nonresidue upstairs, since
\[
 a^{(p^m-1)/r}
   =\left(a^{(p-1)/r}\right)^{1+p+\cdots+p^{m-1}}
   =\left(a^{(p-1)/r}\right)^m\ne1.
\]
The remaining quadratic-even case is proved next. Thus the prior
construction is preserved with an explicit nonvanishing justification.
In the Kummer-cover language, the repair chooses a nonzero vector in the
specified Frobenius eigenspace before passing to its torsor class; a
particular chosen generator need not have a nonzero projection.

\paragraph{CX-ALG-018: the even-degree quadratic case, with no GRH.}
\emph{Statement.} For odd $p$ and even $m$, a supplied
$E=\mathbf F_p[X]/(f)$ of degree $m$ admits a deterministic
$\operatorname{poly}(m\log p)$ construction of a quadratic nonresidue.
This is a fresh complete re-proof of the even-degree clause of inherited
CX\_GRH CX-N13, not a new claim of priority over BIMS17 Corollary 1.3.

\emph{Proof.} Put $Q=p^{m/2}$ and $\alpha=X\bmod f$.
If $p\equiv1\pmod4$ or $4\mid m$, then $Q\equiv1\pmod4$.
The element $u=\alpha^Q-\alpha$ is nonzero, since $\alpha$ has degree
$m$ and is not in the subfield of degree $m/2$. It satisfies $u^Q=-u$.
Consequently
\[
       u^{(p^m-1)/2}=(-1)^{(Q+1)/2}=-1,
\]
so it is the required nonresidue.

It remains that $p\equiv3\pmod4$ and $m\equiv2\pmod4$.
Compute the subfield $K=\mathbf F_{p^2}\subseteq E$ as the kernel of
the $\mathbf F_p$-linear map $z\mapsto z^{p^2}-z$.
Choose $v\in K\setminus\mathbf F_p$ from a computed basis and set
$u=v^p-v\ne0$. Then $u^p=-u$ and $c=u^2\in\mathbf F_p$ is a
nonresidue: a square $c$ would force $u\in\mathbf F_p$.
Because $-c$ is a square, the exponent formula
$t=(-c)^{(p+1)/4}$ gives $t^2=-c$. Thus $i=u/t$ satisfies $i^2=-1$,
and $K=\mathbf F_p(i)$ is explicitly represented.

Start with $z=i$, of order four. Its norm to $\mathbf F_p$ is one.
Whenever the current $z=a+bi$ has norm one, its order is a power of two
at least four; hence $b\ne0$, as the 2-part of $\mathbf F_p^\times$ has
order two. The nonzero numbers
\[
                 \delta_+=(a+1)/2,\qquad \delta_-=(a-1)/2
\]
have product $-b^2/4$, a nonsquare. Exactly one is a square in
$\mathbf F_p$. Determine which by Euler's criterion, choose its square
root $x=\delta^{(p+1)/4}$, and set $y=b/(2x)$.
Using $a^2+b^2=1$ verifies directly that
$(x+yi)^2=a+bi=z$. The norm of this root is $1$ if
$\delta=\delta_+$ and $-1$ if $\delta=\delta_-$.
Replace $z$ by this root, thereby doubling its order, until the norm is
$-1$.

This loop terminates after $O(\log p)$ steps. Indeed the norm-one subgroup
has order $p+1$, whereas the 2-part of $K^\times$ has order
$2^{v_2(p^2-1)}=2^{1+v_2(p+1)}$. An element's order cannot continue
doubling inside the former subgroup. At termination
\[
        z^{(p^2-1)/2}
          =N_{K/\mathbf F_p}(z)^{(p-1)/2}=-1.
\]
The extension $E/K$ has odd degree $m/2$, so this $z$ remains a nonresidue
in $E$. All operations are modular exponentiation or linear algebra in
the supplied $m$-dimensional field, with polynomial bit cost.
\hfill$\square$

The quadratic square-root algorithm then applies in $E$ with this known
nonresidue. In characteristic two the separate square-root problem is
inverse Frobenius, $a\mapsto a^{2^{m-1}}$, and there is no quadratic
nonresidue. These are finite-\'etale and Kummer computations; SGA 1 V
Section 7 and SGA 4 IX 3.2 are the same companions as for CX-ALG-008.

\paragraph{CX-ALG-021: the genus quotient only needs rational Dirichlet-GRH.}
\emph{Statement.} For an imaginary quadratic field $K$, put
$C=\operatorname{Cl}(\mathcal O_K)$ and $G=C/C^2$. The small-prime
spectral and mixing bounds in CX-ALG-020 hold on $G$ under $A_D$, with
$n=2$, $\mathfrak m=(1)$ and $Q=|D_K|$. No Hecke-GRH assumption for
the remaining, non-genus characters of $C$ is needed. If $C$ has exponent
two, this covers the whole class-group graph.
Moreover, the images of prime ideals of prime norm at most
\[
                         (8\log|D_K|+15)^2
\]
generate $C/C^2$ under $A_D$.

\emph{Proof.} Let $H$ be the Hilbert class field of $K$.
By class field theory $\operatorname{Gal}(H/K)=C$. Complex conjugation
preserves $H$: the conjugate of an unramified abelian extension of $K$
is again such an extension, and the maximal one is unique. Its action on
$C$ is inversion, since
$\mathfrak a\overline{\mathfrak a}=(N\mathfrak a)$ is principal.
Actual complex conjugation has order two and restricts nontrivially to
$K$, so
\[
       \operatorname{Gal}(H/\mathbf Q)=C\rtimes\langle c\rangle,
                    \qquad cac^{-1}=a^{-1}.
\]
The subgroup $C^2$ is normal. Hence $F=H^{C^2}$ is Galois over
$\mathbf Q$, and its Galois group is
$(C/C^2)\times\langle c\rangle$, since inversion is trivial on a group
of exponent two. In particular $F/\mathbf Q$ is abelian.

Every character $\chi$ of $C/C^2$ extends to two one-dimensional
characters $\chi_+,\chi_-$ of $\operatorname{Gal}(F/\mathbf Q)$, by
assigning $c$ the values $1,-1$. Directly on the two cosets of
$\operatorname{Gal}(F/K)$,
\[
       \operatorname{Ind}_{\operatorname{Gal}(F/K)}^{
                       \operatorname{Gal}(F/\mathbf Q)}\chi
                         =\chi_+\oplus\chi_-.
\]
Artin induction of L-functions and abelian class field theory over
$\mathbf Q$ therefore express the class-character Hecke function
$L_K(s,\chi)$ as a product of the two associated Dirichlet L-functions.
This includes $\chi=1$, for which the product is $\zeta_K$.
The input $A_D$ gives precisely the zero-location premises for all these
functions.

For the explicit generator bound, fix any nontrivial character $\chi$ of
$C/C^2$. Its kernel in $C$ defines a quadratic unramified extension
$F_\chi/K$ inside $H$. The preceding conjugation argument shows that
$F_\chi/\mathbf Q$ is biquadratic. Its absolute discriminant is
$|D_K|^2$ by the relative-discriminant tower formula and unramifiedness.
Apply Bach--Sorenson Theorem 5.1 to the nonidentity element of
$\operatorname{Gal}(F_\chi/K)$, regarded inside
$\operatorname{Gal}(F_\chi/\mathbf Q)$. Its zeta hypothesis follows from
$A_D$, since this degree-four field is abelian over $\mathbf Q$.
The resulting unramified rational prime satisfies
\[
            p\le(4\log(|D_K|^2)+2.5\cdot4+5)^2
                    =(8\log|D_K|+15)^2.
\]
Its Frobenius is trivial on $K$, so it splits there, and either prime
ideal above it has $\chi$-value $-1$. If the prime ideals in the displayed
range generated a proper subgroup of $C/C^2$, a nontrivial character of
that elementary two-group would be trivial on the subgroup. The prime just
constructed contradicts this. This proves the generator assertion without
paying for the degree of the entire genus field.

The character-by-character proof of JMV Theorem 1.1, author p.~9,
uses only the L-function of the individual character being estimated.
The spectrum of the quotient graph consists exactly of characters
factoring through $C/C^2$ (Remark 1.2(a), p.~4). Thus all its eigenvalue
bounds hold under the premises just established. Lemma 2.1, pp.~7--8,
and the proof of Corollary 1.3, p.~10, give the claimed mixing bound
with $|G|=|C/C^2|$. If $G$ is trivial, the statement is immediate.
\hfill$\square$

This is the genus-level part of the arithmetic Picard action. It sharpens
the input split in CX-ALG-013: the full Hilbert class field may be
nonabelian over $\mathbf Q$, while its genus quotient is abelian and
already covered by $A_D$. It does not give the missing estimates for the
other class characters or a general deterministic class-group algorithm.

\paragraph{CX-GEO-001: existential constructible feasibility.}
\emph{Statement.} Conditional on $M$, deciding whether a quantifier-free
Boolean circuit of rational polynomial equalities and inequations has a
complex solution belongs to $\mathrm{AM}$. Its emptiness problem belongs to
$\mathrm{coAM}\subseteq\Sigma_2^{\mathrm P}$. Polynomial expressions are
given by arithmetic circuits and all original variables are explicitly listed.

\emph{Proof.} Represent each atomic equality $f(x)=0$ by a new Boolean
variable $b$ and a new variable $u$, with equations
\[
 b(b-1)=0,\qquad bf(x)=0,\qquad (1-b)(u f(x)-1)=0.
\]
If $b=1$, these require $f(x)=0$. If $b=0$, they require $f(x)\ne0$
and an inverse $u=f(x)^{-1}$. Conversely any assignment to $x$ extends
with $b$ equal to the truth value of $f(x)=0$; when $b=1$, any $u$ works.
An inequation is the negation of its equality atom.
For each Boolean gate, introduce its output variable $t$ and use
\[
 t-uv=0\quad(\mathrm{AND}),\qquad
 t-u-v+uv=0\quad(\mathrm{OR}),\qquad
 t+u-1=0\quad(\mathrm{NOT}),
\]
where $u,v$ here denote that gate's Boolean inputs. Impose $t(t-1)=0$
for each output if desired, and impose that the final output equals $1$.
Induction through the acyclic Boolean circuit proves that the resulting
polynomial system is feasible exactly when the original formula is true
at some $x$. The construction shares each original arithmetic circuit;
it does not expand a Boolean formula into disjunctive normal form.
Its size and rational-constant bit lengths increase polynomially.
Apply CX-HN-001 to the resulting system and take complements for emptiness.
The reduction concerns equations and inequations over $\mathbf C$, not
ordered-field inequalities. Geometrically it realizes the given constructible
locus as the image of an explicitly presented affine scheme and preserves
nonemptiness.\hfill$\square$

\paragraph{CX-GEO-002: radical membership.}
\emph{Statement.} Conditional on $M$, for
$I=(f_1,\ldots,f_s)\subseteq\mathbf Q[x_1,\ldots,x_n]$ and $g$ given
by rational arithmetic circuits, radical membership $g\in\sqrt I$ is in
$\mathrm{coAM}$ and nonmembership is in $\mathrm{AM}$.

\emph{Proof.} By the Hilbert Nullstellensatz over $\mathbf C$,
$g\notin\sqrt{I\mathbf C[x]}$ exactly when there is a point $x$ with
$f_i(x)=0$ for every $i$ and $g(x)\ne0$. This is equivalent to feasibility of
\[
               f_i(x)=0\quad(1\leq i\leq s),\qquad y g(x)-1=0.
\]
For completeness, the rational and complex membership tests agree. If
$g^m\in I\mathbf C[x]$, a representation by a finite list of polynomial
multiples gives a finite linear system with rational coefficients for those
multipliers' coefficients. A rational linear system having a complex solution
has a rational solution, by row reduction. Thus $g^m\in I$. The converse
follows by extension of scalars. The new system has polynomial size, so
CX-HN-001 gives the asserted classes. In scheme language the question is
whether the principal open $D(g)$ in $\operatorname{Spec}(\mathbf Q[x]/I)$
has a geometric point. Nilpotents matter only through the radical in this
particular problem.\hfill$\square$

\paragraph{CX-GEO-003: reduced containment and equality.}
\emph{Statement.} Conditional on $M$, inclusion and equality of affine
algebraic sets specified by rational arithmetic-circuit equations are in
$\mathrm{coAM}$. This statement does not assert the same bound for
scheme-theoretic ideal inclusion.

\emph{Proof.} Put $X=V(f_1,\ldots,f_s)$ and
$Y=V(g_1,\ldots,g_t)$. Then $X\nsubseteq Y$ exactly when
\[
        f_i(x)=0\quad(1\leq i\leq s),\qquad
        \sum_{j=1}^t u_jg_j(x)=1                              \tag{CX.1}
\]
is feasible. Indeed, the last equation is impossible if all $g_j(x)$ vanish;
if one does not vanish, its inverse supplies a solution with all other
$u_j=0$. For $t=0$ the sum is zero, making (CX.1) impossible, correctly
expressing $Y=\mathbf A^n$. For an empty $f$-list there is no restriction
on $x$, also correctly. Thus noninclusion is in $\mathrm{AM}$ and
inclusion is in $\mathrm{coAM}$.

Equality can be handled by the same single-reduction method. Introduce
$b$ with $b(b-1)=0$. Impose
\[
\begin{aligned}
 b f_i(x)&=0,& (1-b)g_j(x)&=0,\\
 b\left(\sum_j u_jg_j(x)-1\right)&=0,&
 (1-b)\left(\sum_i v_if_i(x)-1\right)&=0.
\end{aligned}
\]
For $b=1$ this is a point of $X\setminus Y$; for $b=0$ it is a point
of $Y\setminus X$. Hence inequality is one HN instance of polynomial size.
Finally, the ideals $(x)$ and $(x^2)$ have the same complex zero set but
are unequal. They show exactly why this reduction tests underlying reduced
sets rather than the structure sheaves or ideal inclusions.
For GCT this test applies when defining equations of the orbit closures
are actually supplied; producing such equations remains a separate gate.
\hfill$\square$

\paragraph{CX-GEO-004: a finite Frobenius test with prescribed excluded primes.}
\emph{Statement.} Let $L/\mathbf Q$ be finite Galois, let $G$ be its Galois
group, and let $E_i=L^{H_i}$ for subgroups $H_i\leq G$ ($i=1,2$).
Let $\chi_i$ be the permutation character of $G/H_i$. Set $n=[L:\mathbf Q]$
and $D=|D_L|$. Suppose the right halves of the critical strips for
$\zeta_L$ and $\zeta_{\mathbf Q}$ are zero-free, as supplied by $M$, or by
$A_D$ if $L/\mathbf Q$ is abelian. If $\chi_1\ne\chi_2$, there is a prime
unramified in $L$ such that
\[
       p\leq B_L=(4\log D+\tfrac52 n+5)^2                  \tag{CX.2}
\]
and the numbers of degree-one prime ideals of $E_1$ and $E_2$ above $p$
differ. Equivalently the numbers of $\mathbf F_p$-points of the two integer
models differ.

For any nonempty finite set $S$ of rational primes put
$d=\prod_{p\in S}p$, $F=\mathbf Q(\sqrt d)$, and $L'=LF$.
The same assertion holds with $p\notin S$ and bound $B_{L'}$, under the
corresponding hypotheses for $\zeta_{L'}$ and $\zeta_{\mathbf Q}$.
If $n'=[L':\mathbf Q]$ and $D'=|D_{L'}|$, an explicit upper estimate is
\[
 n'\leq2n,\qquad \log D'\leq2\log D+n\log|D_F|
                     \leq2\log D+n\log(4d),               \tag{CX.3}
\]
so the entirely explicit larger bound
\[
        \bigl(8\log D+4n\log(4d)+5n+5\bigr)^2            \tag{CX.4}
\]
also suffices. If $S=\varnothing$, use (CX.2) directly.

\emph{Proof.} Choose $\sigma\in G$ such that
$\chi_1(\sigma)\ne\chi_2(\sigma)$. Bach--Sorenson,
\emph{Explicit Bounds for Primes in Residue Classes} (1996), Theorem 5.1,
p.~1729, gives an unramified prime whose Frobenius conjugacy class is that
of $\sigma$, with bound (CX.2). Their Section 2, p.~1720, specifies the
needed zero-free right-half assumption for the top-field zeta function and
$\zeta_{\mathbf Q}$; no Artin-holomorphy conjecture is required.
For an unramified prime the Frobenius orbits on $G/H_i$ correspond to the
prime ideals above $p$, and the orbit lengths are their residue degrees.
Its fixed points therefore count precisely the degree-one primes.
This proves the first assertion.

For the exclusion refinement, every prime of $S$ ramifies in $F$. For odd
$p\mid d$ this follows from the quadratic discriminant formula; if $2\in S$,
then $d$ is even and $D_F=4d$, so 2 ramifies as well. Consequently every
prime in $S$ ramifies in $L'$, since its ramification index in the subextension
$F/\mathbf Q$ divides its index upstairs. Choose any lift of $\sigma$ to
$\operatorname{Gal}(L'/\mathbf Q)$ and apply the same theorem there. The
prime obtained is unramified in $L'$, hence avoids $S$, and its restriction
to $L$ has the required character difference.

We justify (CX.3) without an unstated discriminant estimate. If $F\subseteq L$,
then $L'=L$ and the inequalities are immediate. Otherwise $L$ and $F$ are
linearly disjoint because $F$ is quadratic. If $\omega$ is an integral basis
generator for $\mathcal O_F$, the order $\mathcal O_L[\omega]\subseteq
\mathcal O_{L'}$ has relative discriminant generated by $D_F$. The relative
field discriminant divides the order discriminant. The tower formula gives
$D'\leq D^2|D_F|^n$. Also $n'=2n$ and $|D_F|\leq4d$.
Substitution in (CX.2) yields (CX.4). If $L/\mathbf Q$ is abelian, $LF$ is
also abelian, so the stated $A_D$ alternative survives the exclusion step.
The exceptional source case $L=\mathbf Q$ cannot have unequal characters
for the two subfields in this statement, so causes no gap.\hfill$\square$

\emph{Consequences and limits.} Agreement at all the tested primes implies
$\chi_1=\chi_2$. Hence the corresponding characteristic-zero permutation
representations are isomorphic and their Artin $L$-functions, including
ramified factors formed on inertia invariants, agree. In particular
$\zeta_{E_1}=\zeta_{E_2}$. This is a finite trace certificate, not necessarily
an efficient algorithm in a short presentation of $E_i$: a common Galois
closure may have large degree. Nor is equality of permutation characters
an isomorphism theorem for finite covers, as the next result proves.
The Grothendieck framework is SGA 1 V Section 7 and the Frobenius trace
formalism; cote 98 scans 11--12 are the finite-$G$-set archival companion
carried from the CX\_GRH crosswalk.

\paragraph{CX-GEO-005: equal characters do not identify transitive $G$-sets.}
\emph{Statement.} Let $V=\mathbf F_2^3$, $G=\operatorname{GL}(V)$,
$X=V\setminus\{0\}$ and $Y=V^*\setminus\{0\}$ with the natural and dual
actions. Then $X$ and $Y$ are transitive seven-element $G$-sets with identical
permutation characters but are not isomorphic as $G$-sets.

\emph{Proof.} For $g\in G$, the number of fixed nonzero vectors is
$2^{\dim\ker(g-I)}-1$. The corresponding number of fixed nonzero covectors
is $2^{\dim\ker((g^{-1})^{\mathsf T}-I)}-1$. The matrices
$(g^{-1})^{\mathsf T}-I$ and $g^{\mathsf T}-I$ have the same rank, since
they differ by multiplication by an invertible matrix and sign. Transpose
preserves rank, so both fixed-point counts are equal.

Suppose the two $G$-sets were isomorphic. The stabilizer $H$ of a nonzero
vector, say $e_1$, would stabilize some nonzero covector and therefore its
kernel, a plane $W\subset V$. Such an invariant plane does not exist.
If $e_1\in W$, its image $W/\langle e_1\rangle$ is a line in
$V/\langle e_1\rangle$, invariant under all of
$\operatorname{GL}_2(\mathbf F_2)$ because $H$ induces that full group.
No line is invariant under that group. If $e_1\notin W$, projection identifies
$W$ with the two-dimensional quotient, so $W$ is a complement to
$\langle e_1\rangle$. But $H$ contains all transvections
$(a,v)\mapsto(a+u(v),v)$ with $u$ a linear functional on that quotient.
Choosing $u(v)=1$ for a nonzero projected vector in $W$ moves its lift by
$e_1$, which cannot stay in $W$. Both cases contradict invariance.
Thus no $G$-set isomorphism exists.\hfill$\square$

This is a representation-theoretic counterexample, needing no GRH and no
inverse-Galois assertion. For any connected finite \`etale $G$-torsor that
realizes this group, its two associated connected covers exhibit the same
limitation. No particular unverified realization over $\mathbf Q$ is used.
The example distinguishes the Galois category from its linearized
representation category; the latter is what Frobenius traces recognize.

\paragraph{CX-GEO-006: orbit membership with an explicit action.}
\emph{Statement.} Conditional on $M$, suppose an algebraic group $G$ and its
action on an affine space have polynomial-length rational circuit descriptions,
with every inverted denominator explicitly listed. Given rational points
$v,w$, deciding whether $w\in G(\mathbf C)v$ belongs to $\mathrm{AM}$.

\emph{Proof.} Introduce coordinates for $g$, impose the supplied equations
and inequations defining $G$, and impose the coordinate equations
$g\cdot v=w$. For example, in $\operatorname{GL}_n$ an additional variable
$z$ and equation $z\det(g)-1=0$ enforce invertibility; the determinant has
a polynomial-size arithmetic circuit. Inverted denominators of a rational
action are handled by the same inverse-variable equations. The resulting
constructible-feasibility instance has polynomial size and is equivalent to
orbit membership by definition of the action. Apply CX-GEO-001.
Orbit-closure membership is not represented by this system: points in the
boundary need not equal $g\cdot v$ for any $g$. If closure equations are
provided, CX-HN-001 and CX-GEO-003 apply to those equations; without them
the equation-access problem persists.\hfill$\square$

\paragraph{CX-GEO-007: polynomially parameterized orbit closures already admit a PSPACE test.}
\emph{Statement.} Let the ambient coordinates of $V$ be explicitly listed.
Suppose the action on a supplied rational vector $v$ is given by
division-free arithmetic circuits with rational coefficients for all its
coordinates. In either of the following
cases, membership of a rational point $w$ in the complex orbit closure is
in $\mathrm{PSPACE}$, without either GRH input: (a) the group is
$\operatorname{GL}_n$ and the map $g\mapsto g\cdot v$ extends by the supplied
circuits to a polynomial map on all $\operatorname{Mat}_n$; (b) the group is
$\operatorname{SL}_n$ and the supplied action coordinates are polynomial
in its matrix entries. This statement charges the explicit ambient
coordinate count in the input size.

\emph{Proof.} Guo--Saxena--Sinhababu (2019), Theorem 1.3, p.~5, together
with Lemma 4.1 and Theorem 4.2, pp.~15--16, gives a polynomial-space test
for whether zero lies in the Zariski closure of the image of a polynomial
map specified by arithmetic circuits: this is their approximate polynomial
satisfiability problem. In case (a), $\operatorname{GL}_n$ is dense in
$\operatorname{Mat}_n$. For a polynomial map $F$ on the latter, every closed
set containing $F(\operatorname{GL}_n)$ has closed inverse image containing
a dense subset, hence contains $F(\operatorname{Mat}_n)$. Their image
closures are equal. Apply the cited test to $F-w$.

For case (b), construct a polynomial map from an affine space into
$\operatorname{SL}_n$ whose image contains its Gaussian-elimination big
cell. Products of elementary unipotent matrices parameterize arbitrary
lower and upper unitriangular matrices. A diagonal matrix of determinant
one is a product of adjacent diagonal blocks
$\operatorname{diag}(t,t^{-1})$: for diagonal entries $d_1,\ldots,d_n$,
use $t_i=d_1\cdots d_i$ in positions $(i,i+1)$ for $1\leq i<n$.
Such a block belongs to the image of a
six-factor elementary-unipotent word, because, in its $2$ by $2$ block,
\[
 E_{12}(t)E_{21}(-t^{-1})E_{12}(t)
 E_{12}(-1)E_{21}(1)E_{12}(-1)
       =\operatorname{diag}(t,t^{-1}).
\]
Use independent parameters in each of the six positions of this fixed
word. The resulting map is polynomial and always has determinant one;
the displayed specializations show that it contains every such diagonal
block in its image. Place these words between the lower and upper
unitriangular words. Its image contains every $LDU$ decomposition in the
big cell, which is a dense open subset of $\operatorname{SL}_n$.
The number of elementary factors and parameters is $O(n^2)$ and the
matrix product has polynomial-size circuits. Composing with the supplied
action on $v$ gives a polynomial map with the same image closure as the
orbit map, by the same closed-inverse-image argument. Subtract $w$ and
apply the cited PSPACE algorithm. For $n=1$, the special-linear action is
trivial and membership reduces to $w=v$.\hfill$\square$

This gives an equation-free route for this exact orbit-closure encoding.
It does not give an AM bound from $M$ for approximate satisfiability, and
it does not turn an exponentially long list of ambient coefficients into
a polynomial-length input. Those are separate questions from the theorem.

\paragraph{CX-GEO-008: a finite Frobenius generating certificate.}
\emph{Statement.} Let $L/\mathbf Q$ be finite Galois with group $G$ and
assume the zeta hypotheses of CX-GEO-004. For each unramified rational
prime $p\leq B_L$, choose any representative $g_p$ of its Frobenius class.
The elements $g_p$ generate $G$. More strongly, a subgroup meeting every
Frobenius class observed in this range equals $G$. A sublist of at most
$\lfloor\log_2|G|\rfloor$ representatives suffices to generate $G$.
For a prescribed excluded set $S$, replace the bound by $B_{L'}$ from
CX-GEO-004 and use only primes outside $S$. As before, $A_D$ suffices for
abelian $L/\mathbf Q$, and $M$ supplies the general hypothesis.

\emph{Proof.} Let $H<G$ be proper and let $X=G/H$. Double counting gives
\[
 \sum_{g\in G}|\operatorname{Fix}_X(g)|
      =\sum_{x\in X}|\operatorname{Stab}_G(x)|=|G|.
\]
The identity fixes $|X|>1$ points. If every element fixed at least one,
the left side would exceed $|G|$. Hence there is $g$ with no fixed coset.
Since $g$ fixes $aH$ exactly when $a^{-1}ga\in H$, its conjugacy class
misses $H$. The explicit Chebotarev theorem used in CX-GEO-004 produces
that class at an unramified prime $p\leq B_L$. With excluded primes,
choose a lift to $\operatorname{Gal}(L'/\mathbf Q)$ and apply the same
argument at bound $B_{L'}$; every prime of $S$ ramifies in $L'$ and is
therefore avoided. Thus no proper subgroup can meet every observed class,
and in particular no proper subgroup contains all the chosen $g_p$.

Starting from the trivial subgroup, choose a listed $g_p$ outside the
subgroup generated so far until it equals $G$. Each proper enlargement
has index at least two, so after $k$ selections the generated subgroup has
order at least $2^k$. Consequently $k\leq\lfloor\log_2|G|\rfloor$.
For $G=1$ the empty list suffices.\hfill$\square$

In SGA 1 V Section 7 this is a finite quotient of the arithmetic
fundamental group, with a finite list of Frobenius conjugacy classes
certifying that no proper subgroup contains the image. The bound is
relative to the specified extension and its discriminant; it is not a
uniform presentation of the entire profinite fundamental group. The
archival companion is the finite-$G$-set discussion of cote 98 scans
11--12 recorded in the preceding CX\_GRH crosswalk, and the work target
is ORDER item 30's effective finite Galois-image question.

\paragraph{CX-GEO-009: bounded Laurent degenerations give an AM promise.}
\emph{Statement.} Let $F:\mathbf A^n\to\mathbf A^m$ and
$w\in\mathbf Q^m$ be specified by division-free rational arithmetic
circuits of total size $s$. Supply in unary an upper bound $D\ge1$ for
their syntactic degrees and an integer $B\ge0$.
Conditional on $M$, it is in $\mathrm{AM}$ to decide whether there are
Laurent polynomials
\[
 \beta_i(\varepsilon)=\sum_{j=-B}^{B}c_{i,j}\varepsilon^j,
 \qquad c_{i,j}\in\mathbf C,
\]
such that every coordinate of $F(\beta)-w$ belongs to
$\varepsilon\mathbf C[\varepsilon]$.
Consequently image-closure membership is in $\mathrm{AM}$ on the promise
that every positive instance admits such an arc with $D$ and $B$ bounded
polynomially in the original encoding length. This includes the polynomial
orbit parameterizations of CX-GEO-007 with the same arc promise.

\emph{Proof.} Use the coefficients $c_{i,j}$ as unknowns. At each circuit
gate $g$ of syntactic degree $d_g$, maintain arithmetic circuits for the
coefficients of its substituted Laurent polynomial in the range
$[-Bd_g,Bd_g]$. Input-variable gates have their prescribed coefficient
unknowns; a rational constant has only a constant coefficient. At an
addition or subtraction gate, add or subtract coefficient circuits.
At a multiplication gate, take the usual coefficient convolutions.
Syntactic degrees add at multiplication, so no exponent outside the stated
range can occur. Discard gates unused by the outputs; then each remaining
degree is at most $D$.

There are at most $2BD+1$ coefficient positions per gate, and each product
uses at most $(2BD+1)^2$ pairs. Sharing all earlier coefficient circuits
therefore gives a polynomial-size arithmetic-circuit construction in
$s,n,m,B,D$. For every output impose that all coefficients of exponents
$-BD,\ldots,0$ in $F(\beta)-w$ vanish. These are polynomial equations in
the $n(2B+1)$ unknown coefficients. The resulting HN instance has polynomial
encoding size and is feasible exactly when the displayed bounded Laurent
arc exists. Koiran's theorem gives the claimed AM bound under $M$.
The construction also covers $B=0$, when it is ordinary exact feasibility.

Any such arc puts $w$ in the image closure by GSS Lemma 4.1 and Theorem
4.2, pp.~15--16. Conversely the promise provides a bounded arc for every
point of that closure under consideration, proving the promise corollary.
All oracle queries concern this explicit coefficient system; its bounded
solution field is the exact number-field zeta input. An $A_D$ alternative
requires the bounded abelian-witness promise for this coefficient system,
not merely for the original rational points $v,w$.
\hfill$\square$

The generic GSS bound on Laurent precision can be exponential; writing
that bound in binary would not justify the polynomial-size coefficient
construction. This result therefore isolates a concrete smaller gate for
the promised families: a polynomially bounded degeneration, rather than
access to all equations of the orbit closure. It does not assert that
arbitrary GCT orbit closures meet that promise. The geometric objects are
finite Laurent jets of a parameterized degeneration; the group-action
framework is SGA 3 I 4.6.3.1 and 4.7.1--4.7.2, with cote 28 scan 3 as the
archival companion inherited from CX\_GRH. The work target is ORDER item
28's Koiran/GCT route.

\subsection{Gates retained rather than silently discharged}

\paragraph{CX-LIM-001: general derandomization and GCT.}
The exact remaining premises of the cited routes are uniform hitting sets
or independent circuit-hardness hypotheses. GCT V Theorem 4.7, p.~33,
assumes a specific permanent border-hardness bound. Its Lemma 4.8, on the
same page, extracts a hitting set from a sufficiently explicit separating
system. Neither construction is furnished by a statement about zeros of
Dirichlet or Hecke $L$-functions. This records the boundary of these proofs;
it is not a proof that no logical implication from GRH to general
derandomization could ever exist. No $\mathrm{BPP}=\mathrm P$, integer
factorization in $\mathrm P$, or $\mathrm{VP}\ne\mathrm{VNP}$ is claimed.

\paragraph{CX-LIM-002: positivity and purity.}
The earlier ALPOGE\_GCT P-F7 and P-F10 and PLETH PLE-F5/F7 already show
why pure geometry, nonnegative Hilbert data, efficient counting and actual
orbit-closure multiplicities are different inputs. Their established
counterexamples and corrected positivity-index implications are carried,
not rebranded. The new tests CX-GEO-002--CX-GEO-003 can verify specified
reduced-locus properties once equations are given. They do not construct a
uniform positive multiplicity formula or a strict asymptotic quotient-rank
inequality. Those are the surviving GCT gates.

\paragraph{CX-LIM-003: the separate counting premise.}
The attached CX\_GRH return also assumed an independent theorem denoted
$E$ and drew ETH/counting consequences in CX-E01--CX-E05. Tonight's two
analytic inputs $A_D$ and $M$ do not contain $E$. Those rows therefore do
not become consequences of these two inputs by being present in the same
archive. This is an input classification, not a new assertion about ETH.
